\section{Generating Unconventional Spin--Orbit Torque}

Take the film normal as $\hat{\mathbf z}$ and drive an in-plane electric field
$E_x\hat{\mathbf x}$.  Write $Q_z^{\,a}$ for a spin current flowing toward the
receiving magnetic layer along $z$, with spin polarization along $a$.  The
conventional spin Hall geometry produces $Q_z^{\,y}$; a nonzero $Q_z^{\,z}$
instead supplies an out-of-plane spin polarization, often called a
$\sigma_z$ component.  The spin-current component, its transmission across the
interface, and the torque measured in the receiver are distinct quantities.
The damping-like torque from a transmitted $\sigma_z$ component has the angular
form $\mathbf m\times(\hat{\mathbf z}\times\mathbf m)$, up to the sign of its
amplitude~\cite{ee237-week11,sot-stt-compact-model,zhang2024revisiting}.
This torque vanishes at exact collinearity, $\mathbf m=\pm\hat{\mathbf z}$;
the subsequent switching dynamics must therefore be evaluated for the actual
initial state and other torques~\cite{zhang2024revisiting}.
The examples below describe ways to \emph{allow or generate} this component;
symmetry alone does not set its magnitude or establish that switching occurs.

\subsection{Lower the Crystal Symmetry}

A low-symmetry spin source can permit an out-of-plane torque for one in-plane
current direction and forbid it for another.  For example, with a surviving
$yz$ mirror, $E_x$ and the axial spin component $s_z$ both reverse, whereas
$E_y$ does not.  A $z$-directed flow is unchanged by that mirror, so
$Q_z^{\,z}\propto E_x$ is allowed while $Q_z^{\,z}\propto E_y$ is forbidden.
WTe$_2$/ferromagnet bilayers exhibit this current-direction-dependent
out-of-plane torque; the crystal orientation therefore matters when specifying
the device axes~\cite{ee237-week11,macneill2017wte2}.

\subsection{Use Magnetic Order to Rotate the Spin Polarization}

The magnetization of a spin-source ferromagnet can remove symmetries that
forbid spin-current components in a nonmagnetic source.  Davidson \emph{et al.}
consider $E_x$ and the operations $\sigma_{xz}$, $\sigma_{xy}$, and $C_2^x$.
With no magnetization, only $Q_z^{\,y}$ among the three $z$-flowing spin
components survives all three operations.  For a source magnetization
$\mathbf m_s\parallel\hat{\mathbf x}$, only $C_2^x$ remains: it reverses both
the $z$ flow and the $z$ spin, making $Q_z^{\,z}$ symmetry-allowed.  The
spin-rotated term in the perspective can be written schematically as
\begin{equation}
    \mathbf Q_z^{\,R}\propto
    \mathbf m_s\times(\hat{\mathbf z}\times\mathbf E),
    \qquad
    \mathbf m_s\parallel\hat{\mathbf x},\quad
    \mathbf E\parallel\hat{\mathbf x}
    \ \Longrightarrow\ Q_z^{\,z}\ne 0\ \text{is allowed}.
    \label{eq:ferromagnetic-spin-rotated-current}
\end{equation}
Reversing $\mathbf m_s$ reverses this term.  This is a statement about the
source's allowed spin-current polarization, not a prediction of its size or of
the receiver's switching outcome~\cite{davidson2020spin-currents}.

\paragraph{Source status.}
\textbf{Reference needed.}  The detailed ferromagnetic selection rule and
Eq.~\eqref{eq:ferromagnetic-spin-rotated-current} come from the Davidson
perspective and its processed explanation; an approved project reference is
still needed for these details.

A CoFeB/Ti/CoFeB trilayer provides a device example: the bottom in-plane
magnetized layer or its interface generates both in-plane and out-of-plane
spin-polarization components, and the upper perpendicular CoFeB layer is the
receiver.  Reversing the bottom-layer magnetization reverses the reported
switching polarity~\cite{yang2024ferromagnetic-trilayer}.

\paragraph{Source status.}
\textbf{Reference needed.}  The trilayer result has no approved project
reference yet.

\subsection{Design Electrical Asymmetry in a Polycrystalline Source}

Broken electrical symmetry can also produce a $z$-polarized spin current
without relying on a low-symmetry single crystal.  Liu \emph{et al.} report
$z$-spin damping-like and field-like torques in sputtered polycrystalline
heavy-metal devices with transverse and perpendicular electrical asymmetry.
They vary the device length, width, thickness, and metal to tune the response,
and demonstrate nearly full field-free switching of perpendicular FeCoB.
The device geometry is therefore part of this generation
route~\cite{liu2024electric-asymmetry}.

\paragraph{Source status.}
\textbf{Reference needed.}  The electrical-asymmetry mechanism and device
result have no approved project reference yet.

In each route, report the in-plane current direction, film normal, source
magnetization where present, and receiving layer before assigning a sign to
$\sigma_z$ or to its torque amplitude~\cite{ee237-week11,
    sot-stt-compact-model}.
